\sectionbreak \section*{
	\gostTitleFont
	\redline
	ЗАКЛЮЧЕНИЕ
}

\subtitlespace

{\gostFont
	\par \redline В рамках данного дипломного проекта была разработана серверная часть приложения заказа такси на основе микросервисной архитектуры с использованием языка Java и экосистемы Spring Boot.
	
    \par \redline В процессе выполнения дипломного проекта были решены следующие задачи:
	
	\begin{itemize}[leftmargin=2.15cm, labelwidth=0.65cm, labelsep=0.0cm] 
		
		\item[\theitemcntr.] Проведён анализ предметной области сервисов заказа такси, выделены основные бизнес-сущности и процессы.		
		\addtocounter{itemcntr}{1}
		
		\item[\theitemcntr.] Выполнен сравнительный анализ монолитной и микросервисной архитектур, обоснован выбор микросервисного подхода.		
		\addtocounter{itemcntr}{1}
		
		\item[\theitemcntr.] Спроектирована архитектура системы: схема взаимодействия микросервисов, структура БД, REST API конечных точек.		
		\addtocounter{itemcntr}{1}
		
		\item[\theitemcntr.] Реализованы ключевые модули: сервис поездок с конечным автоматом, геолокация с кэшированием, авторизация через JWT, асинхронное взаимодействие через Kafka.		
		\addtocounter{itemcntr}{1}
		
		\item[\theitemcntr.] Разработаны модульные и интеграционные тесты с использованием TestContainers, Cucumber и REST Assured.		
		\addtocounter{itemcntr}{1}
		
		\item[\theitemcntr.] Выполнена контейнеризация системы с помощью Docker, настроен мониторинг на базе Prometheus, Grafana, Tempo и Loki.		
		\addtocounter{itemcntr}{1}
		
		\item[\theitemcntr.] Проведён расчёт технико-экономической эффективности, определена себестоимость и отпускная цена продукта.		
		\addtocounter{itemcntr}{1}

		\setcounter{itemcntr}{1}
	\end{itemize} 

	\par \redline Разработанная система удовлетворяет следующим ключевым требованиям:
	
	\begin{itemize}[leftmargin=2.15cm, labelwidth=0.65cm, labelsep=0.0cm] 
		
		\item[\theitemcntr.] Микросервисная архитектура с независимым хранением данных для каждого сервиса.	
		\addtocounter{itemcntr}{1}
		
		\item[\theitemcntr.] Синхронное взаимодействие через REST API и асинхронное - через брокер сообщений Kafka.		
		\addtocounter{itemcntr}{1}
		
		\item[\theitemcntr.] Аутентификация и авторизация на основе JWT с ролевой моделью доступа.		
		\addtocounter{itemcntr}{1}
		
		\item[\theitemcntr.] Кэширование часто запрашиваемых данных с использованием Redis.		
		\addtocounter{itemcntr}{1}
		
		\item[\theitemcntr.] Контейнеризация и мониторинг для обеспечения отказоустойчивости и наблюдаемости.		
		\addtocounter{itemcntr}{1}
		
		\setcounter{itemcntr}{1}
	\end{itemize}  
	 
	\par \redline Проведённое тестирование подтвердило корректность реализации ключевых бизнес-сценариев: регистрация пользователей, создание заказа, назначение водителя, обработка переходов между статусами поездки и обмен событиями через Kafka. 

	\par \redline Для разработанного программного обеспечения выполнен расчёт экономической эффективности: полная себестоимость составила 5020,57 белорусских рублей, отпускная цена - 7410,36 белорусских рубля, чистая прибыль - 923,78 белорусских рубля. Полученные значения подтверждают экономическую целесообразность разработки.
	
}

\setcounter{subchaptercntr}{1}
\setcounter{formulacntr}{1}
\setcounter{imagecntr}{1}
